\chapter{Reproducibilidad}

Este anexo reúne lo necesario para repetir la construcción de los conjuntos, el
entrenamiento de los modelos y la verificación del sistema. No es un manual de
instalación: es la lista de condiciones y comandos que determinan los resultados
reportados.

\section{Entorno de ejecución}

El entorno se declara con archivo de bloqueo de dependencias, de modo que las
versiones efectivamente usadas quedan registradas y no dependen de la fecha de
instalación. La Tabla~\ref{tab:anexo_entorno} resume el entorno de entrenamiento y
el de operación.

\begin{longtable}{
        P{0.23}
        P{0.43}
        P{0.25}
    }
    \caption{Entorno de entrenamiento y de operación}
    \label{tab:anexo_entorno}\\

    \toprule
    \rowcolor{gray!20}
    \textbf{Componente} &
    \textbf{Versión y condición} &
    \textbf{Uso} \\
    \midrule
    \endfirsthead

    \multicolumn{3}{c}{\textit{Continuación de la Tabla
    \ref{tab:anexo_entorno}}} \\
    \toprule
    \rowcolor{gray!20}
    \textbf{Componente} &
    \textbf{Versión y condición} &
    \textbf{Uso} \\
    \midrule
    \endhead

    Python & $3{,}12$, con gestor de dependencias y archivo de bloqueo &
    Servicios, trabajador y herramientas \\
    \midrule

    PyTorch y torchvision & $\geq 2{,}4$ y $\geq 0{,}19$ &
    Entrenamiento e inferencia \\
    \midrule

    Ultralytics & $\geq 8{,}4$ & Detección de vehículos y de placas \\
    \midrule

    OpenCV & $\geq 4{,}10$, sin interfaz gráfica &
    Captura, preprocesamiento y calibración \\
    \midrule

    Acelerador & GPU de $4$ GB de memoria de video &
    Entrenamiento y compilación de motores \\
    \midrule

    Almacenamiento & NVMe para conjuntos, recortes y pesos &
    Fuente de verdad del material \\
    \midrule

    Servicios de operación & PostgreSQL $16$ y Redis $7$ en contenedores &
    Persistencia y cola de tareas \\
    \midrule

    Compilación de modelos & ONNX, ONNX Runtime y TensorRT &
    Motor especializado por dispositivo \\
    \bottomrule
    \multicolumn{3}{r}{\footnotesize\textit{Fuente: Elaboración propia}} \\*
\end{longtable}

\section{Semillas y determinismo}

La partición de ablación del reconocimiento de matrículas se fija con semilla
$7$ y proporción de prueba del $20\%$, y el archivo de protocolo generado registra
la semilla y los tamaños resultantes. La generación de placas sintéticas usa la
misma semilla, de modo que el conjunto sintético es reproducible byte a byte
mientras no cambie el generador.

Los conjuntos de tipo y color no se remuestrean: se derivan del pool congelado de
recortes, y la partición se establece por sesión de grabación. Cambiar la semilla
en esos conjuntos no cambia los resultados, porque la partición no es aleatoria por
imagen.

\section{Comandos por etapa}

Los comandos se ejecutan desde el directorio del software. Cada bloque produce un
artefacto verificable.

\subsection{Construcción y verificación de los conjuntos}

\begin{verbatim}
# Reconstruye los conjuntos de tipo y color desde el pool.
FORCE=1 ./train_all.sh datasets

# Verifica clases y tamaños antes de entrenar.
for d in "type dataset v3" "color dataset v5"; do
  grep '^nc:' "$DATASETS/$d/data.yaml"
  for s in train valid test; do
    ls "$DATASETS/$d/$s/images" | wc -l
  done
done
\end{verbatim}

El valor esperado es $5$ clases con $2.783$, $590$ y $633$ muestras para tipo, y
$11$ clases con $2.482$, $539$ y $551$ para color. La verificación es obligatoria
porque el script de construcción omite un directorio existente sin comprobar su
integridad: un conjunto truncado por una copia fallida se saltearía en silencio.

\subsection{Ablación y entrenamiento del reconocimiento}

\begin{verbatim}
uv run python tools/prepare_ocr_ablation.py --dataset "$DATASETS"
uv run python tools/generate_synth_plates.py \
    --dataset "$DATASETS" --count 360000 --seed 7

uv run python -u tools/train_ocr_crnn.py \
    --dataset "$DATASETS" \
    --stage stage1_synth --epochs-synth 8 \
    --lr-schedule fixed --save-dir "$RUNS/ocr_v2"

uv run python -u tools/train_ocr_crnn.py \
    --dataset "$DATASETS" \
    --stage stage2_real --epochs-real 500 \
    --lr-real 1e-3 --lr-schedule fixed --synth-mix-ratio 0 \
    --save-dir "$RUNS/ocr_v2"
\end{verbatim}

Las banderas no son arbitrarias: cada una reemplaza una configuración que se midió
y se descartó. La tasa de aprendizaje fija evita el colapso del decodificador hacia
la secuencia vacía que produjo el planificador por meseta, y la ausencia de mezcla
sintética en la segunda etapa mantiene el ajuste fino sobre material real.

\subsection{Clasificadores, detector de placas y evaluación}

\begin{verbatim}
./train_all.sh type       # clasificador de tipo
./train_all.sh color      # clasificador de color
./train_all.sh platedet   # detector de placas
./train_all.sh all        # todas las etapas en orden

uv run python tools/eval_type_classifier.py \
    --dataset "$DATASETS/color dataset v5" \
    --checkpoint "$DATASETS/runs/resnet18_final_color_v5_best.pt" \
    --split test

uv run python tools/report_metrics.py \
    --metrics "$DATASETS"/runs/*_metrics.json \
    --title "Clasificadores" --out ../reports/clasificadores.md
\end{verbatim}

Cada checkpoint se acompaña de su archivo de métricas. El archivo no es un adorno:
el motor de inferencia lee de allí el orden de las clases, de modo que un peso sin
su archivo de métricas no puede servirse y el trabajador se niega a arrancar.

\section{Artefactos}

La Tabla~\ref{tab:anexo_artefactos} enumera los artefactos que un tercero necesita
para auditar los resultados, con su ubicación relativa al directorio del material.

\begin{table}[H]
    \centering
    \caption{Artefactos y contenido}
    \label{tab:anexo_artefactos}
    \begin{tabular}{ll}
        \toprule
        \textbf{Artefacto} & \textbf{Contenido} \\
        \midrule
        \texttt{pool\_20261021} & Recortes de vehículo congelados, $5$ clases \\
        \texttt{type dataset v3} & Partición de tipo con archivo de clases \\
        \texttt{color dataset v5} & Partición de color con archivo de clases \\
        \texttt{plates/} & Placas reales y sintéticas, manifiesto \\
        \texttt{plates\_det/} & Cajas propuestas, verificadas y texto \\
        \texttt{runs/} & Pesos de los clasificadores y métricas \\
        \texttt{runs/ocr\_v2/} & Etapas del reconocedor y métricas finales \\
        \texttt{ablation/protocol.json} & Semilla y tamaños de la ablación \\
        \bottomrule
    \end{tabular}
\end{table}

\section{Verificación del sistema}

El software del sistema de gestión se verifica con su propio banco de pruebas. La
ejecución completa comprende las pruebas unitarias y las de integración; estas
últimas requieren un servidor PostgreSQL accesible y se omiten cuando no lo hay,
de modo que la suite puede correr en un equipo sin servicios levantados.

\begin{verbatim}
uv run pytest            # unitarias; omite integración sin base
uv run pytest -m integration       # requiere PostgreSQL accesible
uv run alembic -c backend/alembic.ini upgrade head
\end{verbatim}

El resultado reportado en este trabajo es de $140$ pruebas unitarias y $4$ de
integración aprobadas, con una cobertura del $88\%$ sobre el código de la
aplicación. Las pruebas de integración aplican las migraciones sobre una base
vacía, ejercitan el registro y el cierre de una sesión y verifican el efecto en la
base de datos.

\section{Condiciones que limitan la reproducción}

\begin{itemize}
    \item \textbf{Los motores compilados son específicos del dispositivo.} Un
        motor generado para un acelerador y un controlador determinados no es
        portable a otro equipo, por lo que cada destino requiere su propia
        compilación. Los resultados de latencia no son trasladables entre equipos.
    \item \textbf{El material de origen no se distribuye con el documento.} Las
        imágenes provienen de las cámaras del recinto y su publicación está sujeta
        a las condiciones de tratamiento de datos descritas en el capítulo de
        diseño. Lo reproducible sin ellas son los procedimientos, las particiones
        y las métricas.
    \item \textbf{Una corrida no es un resultado.} La varianza medida entre
        corridas idénticas es de $1$ a $3{,}4$ puntos de F1 macro, de modo que una
        diferencia menor a ese rango no debe interpretarse como mejora. Reproducir
        un valor puntual exige repetir la corrida, no solo el comando.
    \item \textbf{El episodio de integridad del almacenamiento.} El disco externo
        de trabajo devolvió errores de entrada y salida en una fracción de los
        recortes y truncó archivos durante la transferencia. El material se migró
        a almacenamiento NVMe, que pasó a ser la única fuente de verdad; cualquier
        reproducción debe partir de esa copia y no del disco original.
\end{itemize}
